\section{Decapod: A Reference Governance Kernel}
\label{sec:system}

To validate this model, we present the prototype implementation of \textbf{Decapod}, a local-first, repo-native governance kernel. Decapod does not serve as an agent wrapper or model proxy; rather, it sits as a control plane between the agent's runtime process and the target repository.

\subsection{System Architecture}
Decapod operates entirely on the local file system within the target repository under a hidden configuration directory, \texttt{.decapod/}. This local-first structure ensures that all logs, trajectories, and parameters remain developer-controlled, avoiding dependency on remote orchestration platforms.

In our reference design, the core execution flow is managed by three sub-modules:
\begin{enumerate}
    \item \textbf{Workspace Custodian}: Spawns isolated execution environments. When a task is initialized, the Custodian uses Git worktrees to create an ephemeral directory slice linked to a task-specific branch. It is designed to configure shell path controls, filtering out unauthorized system commands and blocking access to environment variables not explicitly permitted by the custody configuration.
    \item \textbf{Trajectory Logger}: Intercepts standard input, output, and file changes. In our prototype, shell commands invoked by the agent run through a pseudo-terminal (PTY) wrapper. The logger plans to hash workspace files before and after each command to generate accurate diff mappings ($\Delta(W_i)$) and append this data to the run log.
    \item \textbf{Validation Gatekeeper}: Enforces the test suites defined in the intent manifest. The Gatekeeper is intended to intercept completion requests, execute the validation suite within the isolated custody context, and compile the outcomes.
\end{enumerate}

\subsection{Intent Manifest and Workflows}
The lifecycle of a Decapod run is guided by a structured JSON manifest, \texttt{.decapod/intent.json}, containing:
\begin{verbatim}
{
  "intent_id": "task-07-fix-auth",
  "allowed_scopes": ["src/auth/", "tests/"],
  "allowed_tools": ["python", "pytest", "git-diff"],
  "validation_gates": [
    {
      "name": "Unit Tests",
      "command": "pytest tests/test_auth.py"
    },
    {
      "name": "Linter",
      "command": "flake8 src/auth/"
    }
  ]
}
\end{verbatim}

The agent is provided with this manifest at startup. It communicates with the kernel using simple file markers. When the agent asserts completion, the Gatekeeper is designed to intercept this signal, lock the workspace from further agent writes, and execute the validation gates.

\subsection{Cryptographic Proof Compilation}
If all validations pass, the prototype kernel compiles the final \texttt{proof.json} file. This file is designed to contain:
\begin{itemize}
    \item The target repository commit hash.
    \item The base64-encoded output of the successful validation commands.
    \item The final file modifications (diff).
    \item A cryptographic hash of the compiled trajectory.
    \item A local signature generated using a developer-configured GPG or SSH key.
\end{itemize}
This proof artifact is committed directly to the task branch, allowing CI/CD systems or human reviewers to verify the integrity of the agent's work without re-running the entire execution.
